\scanpage{8}
3. Supposons
\[
 \boxed{B=\prod_{i\in I}B_i}. \tag*{(3.1)}
\]
Donc $f$ est donné par une famille de
\[
 f_i:A\longrightarrow B_i \tag*{(3.2)}
\]
et $g$ est connu quand on connaît les adjoints
\[
 g_i:B_i\longrightarrow A, \tag*{(3.3)}
\]
par
\[
 g\bigl((X_i)_{i\in I}\bigr)=\prod_{i\in I}g_i(X_i) \tag*{(3.4)}
\]
Le composé $\varphi=fg$ est alors un \unsure{hom|hom.}
\[
 \varphi:\prod B_i\longrightarrow\prod B_i
\]
et est connu quand on connaît les
\[
 \mathrm{pr}_{j}\,\varphi:\prod_i B_i\mathrel{\unsure{-|\to}}B_j
\]
On trouve
\[
 \mathrm{pr}_j\,\varphi((X_i))=f_j\left(\prod_{i\in I}g_i(X_i)\right). \tag*{(3.5)}
\]
Soit
\[
 \varphi_{ji}\overset{\text{déf}}{=}f_jg_i\deleted{\unsure{g_j f_i}}:\ B_i\longrightarrow B_j \tag*{(3.6)}
\]
et supposons que les $f_j$ commutent aux produits de type $I$ (p. ex. $f_j$ exacts à g. et $I$ fini). On a alors
\[
 \mathrm{pr}_j\,\varphi((X_i))\simeq\prod_{i\in I}\varphi_{ji}(X_i) \tag*{(3.6)}
\]
et les données de $\alpha,\lambda$ reviennent à celles de
\[
 \prod_i\varphi_{ki}(X_i)\xrightarrow{\bar\alpha_{\deleted{\unsure{j}}k}}X_{\deleted{\unsure{j}}k} \tag*{(3.7)}
\]
\[
 \prod_i\varphi_{\deleted{\unsure{j}}ki}(X_i)\xrightarrow{\bar\lambda_k}
 \prod_{\deleted{\unsure{j}}i,j}\varphi_{kj}\varphi_{ji}(X_i) \tag*{(3.8)}
\]
i.e.
\[
 \prod_\alpha\varphi_{k\alpha}(X_\alpha)\xrightarrow{\bar\lambda_{k,j,i}}\varphi_{kj}\varphi_{ji}(X_i)
\]
% archive sheet number: 7
